\documentclass[10pt,letterpaper]{article}
\usepackage{cornell-notes}

\title{Annex C.02: 02-use-with-iso-iec-ieee-42020}
\author{}
\date{}
\setDocTitle{Annex C.02: 02-use-with-iso-iec-ieee-42020}
\setDocSubtitle{ISO/IEC/IEEE 42010:2022}
\setDocOwner{ISO 42010 Cornell Notes}
\setDocDate{\today}
\setCornellCollection{ISO/IEC/IEEE 42010:2022 Cornell Notes}
\setCornellUnitType{Annex}
\setCornellUnitNumber{C.02}
\setCornellUnitTitle{02-use-with-iso-iec-ieee-42020}
\hypersetup{pdftitle={Annex C.02: 02-use-with-iso-iec-ieee-42020},pdfsubject={Cornell notes},pdfauthor={}}

\begin{document}
\maketitle
\begin{CornellOverviewBox}{Overview}
{\Large\bfseries\textcolor{CNavy}{Annex C.2: Use with ISO/IEC/IEEE 42020}}\\[3pt]
{\small\textbf{Classification:} Informative annex}\\
{\small\textbf{Learning objective:} Use the informative guidance on use with iso/iec/ieee 42020 to reinforce the normative and conceptual material without treating the guidance itself as mandatory.}
\end{CornellOverviewBox}
\vspace{3pt}

\begin{CornellNotesTable}
\CornellNoteRow{Purpose / key concept}{ISO/IEC/IEEE 42020 specifies requirements, recommendations and permissions for architecture processes suitable for the enterprise, organization or project.}
\CornellNoteRow{Core details}{\begin{itemize}[leftmargin=*,nosep,topsep=1pt]
\item In this reference document, AD is mainly addressed by a process called “architecture elaboration”, explaining how to describe or document an architecture in a sufficiently complete and correct manner for the intended uses of the architecture.
\item Nevertheless, another ISO/IEC/IEEE 42020 process, called “architecture conceptualization”, can be undertaken to serve as the basis for an AD in order to characterize the problem space and determine suitable solutions that address stakeholder concerns, achieve architecture objectives and meet relevant requirements.
\item Both architecture conceptualization and architecture elaboration are specified in ISO/IEC/IEEE 42020 through an extensive set of activities, tasks, outcomes and work products.
\item Recognizing that particular projects or organizations do not need to use all of the processes specified by This section, full conformance and tailored conformance are defined to accommodate a flexible implementation approach to claim conformance to ISO/IEC/IEEE 42020.
\end{itemize}}
\CornellNoteRow{Active recall}{\begin{itemize}[leftmargin=*,nosep,topsep=1pt]
\item What is the central role of Use with ISO/IEC/IEEE 42020?
\item How does Use with ISO/IEC/IEEE 42020 connect to adjacent architecture-description concepts?
\end{itemize}}
\end{CornellNotesTable}


\begin{CornellSummaryBox}{Summary}
ISO/IEC/IEEE 42020 specifies requirements, recommendations and permissions for architecture processes suitable for the enterprise, organization or project. This material is informative guidance and does not itself create a conformance requirement.
\end{CornellSummaryBox}

\end{document}
